\section{Model Deployment and Application}
\label{sec:deployment}

The proof of concept is a Streamlit application deployed on Streamlit
Community Cloud at \url{https://team10-it5006-olist.streamlit.app}, built from
the \code{deployment/} folder of the repository. It serves all three deployed
models, and because the Olist extract ends in 2018 it does what the brief
anticipates: it predicts on held-out historical records that stand in for new
orders, and it can replay the test window a day at a time.

\subsection{Architecture}

\begin{table}[htbp]
\centering
\caption{The deployed artifacts, as pickled. Thresholds and the margin-to-risk
mapping were fitted on the validation period; nothing was refitted for
deployment, so the served models are exactly the ones Section~\ref{sec:results}
evaluates.}
\label{tab:deployed}
\footnotesize
\begin{tabular}{@{}p{2.3cm}p{5.6cm}p{2.0cm}p{4.3cm}@{}}
\toprule
\textbf{Problem} & \textbf{Estimator served} & \textbf{Artifact} & \textbf{Decision rule} \\
\midrule
1A lead time & \code{VotingRegressor}: LinearRegression + HistGradientBoosting, weights 1:2 & 381\,KB & re-quote with the estimate \\[2pt]
1B late flag & logistic map of (predicted $-$ promised days) to a probability & 1\,KB & flag at $p_{\text{late}} \ge 0.11$ \\[2pt]
2 low review & \code{StackingClassifier}: $\ell_2$ LogisticRegression + HistGradientBoosting, logistic meta-learner & 321\,KB & flag at $p_{\text{low}} \ge 0.19$; rank for the top-$K$ \\
\bottomrule
\end{tabular}
\end{table}

Nothing in the application reimplements the training-time feature engineering.
A form submission or an uploaded row is turned back into an order frame and
pushed through the same \code{src/features.py} builders the models were fitted
on, which is what stops a feature being computed one way in training and
another way in production \citep{sculley2015debt}. The one thing a form cannot
supply is a seller's track record, so \code{models/history\_lookup.csv} carries
those aggregates at each level of the same fallback chain the training code
uses. The artifacts total under 1\,MB and record the \code{scikit-learn}
version that wrote them; \code{requirements.txt} pins that version, the loader
raises a visible warning if it ever drifts, and the app runs on Python 3.12,
the interpreter the pinned numpy publishes wheels for. Models are loaded
lazily, so a visitor who only browses the exploratory tabs never pays for
unpickling one.

\subsection{What the application does}

The application has eight tabs (Figures~\ref{fig:app-predict} to
\ref{fig:app-card}). Five are the exploratory analysis of Section~\ref{sec:eda}
with sidebar filters. \textbf{Predict} scores a single order entered as it
would look at checkout and returns the expected lead time with its typical
error, the risk of missing the promised date against the tuned threshold, and,
conditional on that forecast, the low-review risk; it also scores an uploaded
CSV of up to 5{,}000 orders and builds the \textbf{triage list}: rows ranked by
review risk with the top-$K$ marked for a capacity the user enters, and the
precision and recall that capacity achieved on the held-out window beside it.
\textbf{Replay a day} is new in this phase (Figure~\ref{fig:app-replay}): the
user picks any purchase day in the test window and a daily capacity, sees what
the deployed models predicted for every order placed that day, ranked the way
the customer-experience team would work them, and a toggle reveals what
actually happened --- which orders left a 1--2 star review, how long delivery
really took --- with the precision of the top-$K$ on that day and on every
other day of the window. For 16 July 2018 at ten contacts, eight of the ten
would have reached a would-be low reviewer. \textbf{Model card} states what is
deployed as read from the artifacts (Table~\ref{tab:deployed}), the training
protocol, every searched model's selected hyperparameters, the fold-to-fold
variation of every candidate, the Precision@$K$ table, the situations in which
the models should not be trusted \citep{mitchell2019modelcards}, and a
calculator for the pilot of Section~\ref{sec:pilot}. Every number it shows is
read from the result tables this report quotes; the application recomputes
nothing.

\subsection{Robustness, reachability and reproducibility}

Invalid input produces a specific message rather than a traceback: a promised
date before the purchase date, a non-positive price, more sellers than items, a
malformed, empty or oversized upload, a missing column, an unknown category or
state, a missing artifact, a version mismatch. The data loader asserts the row
count of all nine source tables and stops the app if they change. There are no
API endpoints and no credentials anywhere: the batch CSV contract, documented
with an example in \code{deployment/README.md}, is the programmatic surface.

The Phase 1 dashboard became inaccessible when the GitHub account it was
deployed from was suspended, which took the Community Cloud app down with it.
The application now deploys from the team's current repository, and two things
keep it reachable, both recorded in the README: Community Cloud puts an app
without visitors for twelve hours to sleep behind a one-click wake-up button
that any visitor can press, and the Python version must stay at 3.12 in the
app's settings. A fresh clone reproduces everything, from the committed source
tables to every table, figure and number in this document
(Appendix~\ref{app:repro}).
